\subsection{Algoritmos de optimización}
\label{sec:algoritmos_de_optimizacion}

Los algoritmos basados en gradientes se pueden clasificar principalmente por la cantidad de información que utilizan. La familia más común son los métodos de primer orden, como el algoritmo de gradiente descendente (GD) \cite{Cauchy1847}. Estos consisten en actualizar en cada iteración los parámetros moviéndose en la dirección opuesta al gradiente (la dirección de máximo decrecimiento) de la función de costo. Este fue, de hecho, el método utilizado para minimizar la log-verosimilitud del ejemplo anterior \eqref{eq:log_likelihood} y ajustar los parámetros del modelo.

Existen también métodos de segundo orden, como el método de Newton-Raphson \cite{Nocedal2006}, que no solo utilizan el gradiente (primera derivada), sino también la matriz Hessiana (segunda derivada) para modelar la curvatura de la función. Aunque estos métodos pueden converger más rápido, el costo de calcular la Hessiana es computacionalmente prohibitivo para la mayoría de los problemas modernos \cite{Goodfellow2016}. Por esta razón, los métodos de primer orden como el algoritmo de gradiente descendente son los utilizados en el entrenamiento de redes neuronales.

El método de gradiente descendiente consiste en actualizar los parámetros en la dirección contraria al gradiente de la función de costo ($\nabla_{\theta} L(y, \hat{y})$), es decir, en la dirección que reduce el error \cite{Cauchy1847}, modulado por una tasa de aprendizaje $\eta$. Esta tasa es un hiperparámetro, que determina el tamaño del paso en cada iteración.

\begin{equation}
    \label{eq:gradient_descent}
    \theta_{n+1} = \theta_n - \eta \nabla_{\theta} L(y, \hat{y})
\end{equation}

El gradiente de una función define la dirección de mayor incremento de la misma en un punto determinado. Por lo tanto, al restar el gradiente en la actualización de los parámetros, se avanza en la dirección de mayor decremento de la función de costo, buscando minimizarla iterativamente.

Es posible también caer en mínimos locales de la función de costo, dependiendo de la forma de la función y del punto de inicio. Sin embargo, en muchos casos prácticos, el gradiente descendiente puede encontrar soluciones suficientemente buenas, incluso si no son óptimas globales \cite{Bishop2006}.

En esta búsqueda iterativa, la elección de la tasa de aprendizaje $\eta$ es crucial. Una tasa de aprendizaje demasiado alta puede provocar que el algoritmo diverja, mientras que una tasa demasiado baja puede resultar en una convergencia lenta \cite{Bishop2006}.

Una variante de GD es el descenso por gradiente estocástico (SGD, por sus siglas del inglés \textit{Stochastic Gradient Descent}) \cite{RobbinsMonro1951}. La diferencia con GD es que en lugar de calcular el gradiente utilizando todo el conjunto de datos $\mathcal{D}$, como en el caso del gradiente descendente tradicional, divide $\mathcal{D}$ en pequeños subconjuntos llamados \textit{minibatches} y calcula el gradiente utilizando solo un \textit{minibatch} en cada iteración. A su vez, los \textit{minibatches} son seleccionados de forma aleatoria en cada iteración, lo que introduce ruido en el proceso de optimización, mejorando la capacidad del algoritmo para escapar de mínimos locales \cite{Goodfellow2016}.

Otra variante ampliamente utilizada es el método de ADAM (\textit{Adaptive Moment Estimation}) \cite{kingma2017adammethodstochasticoptimization}. Este algoritmo combina las ventajas del descenso por gradiente estocástico con la adaptación de la tasa de aprendizaje para cada parámetro, utilizando estimaciones de los primeros y segundos momentos del gradiente. ADAM ha demostrado ser eficaz en una amplia variedad de problemas de aprendizaje automático, especialmente en el entrenamiento de redes neuronales profundas \cite{Goodfellow2016} como se verá en la sección \ref{sec:entrenamient_redes_neuronales}.
